
\documentclass{article}
\usepackage{colortbl}
\usepackage{makecell}
\usepackage{multirow}
\usepackage{supertabular}

\begin{document}

\newcounter{utterance}

\centering \large Interaction Transcript for game `cladder', experiment `full\_v1.5\_default', episode 188 with reflexion{-}llama{-}3.1{-}8b{-}instruct.
\vspace{24pt}

{ \footnotesize  \setcounter{utterance}{1}
\setlength{\tabcolsep}{0pt}
\begin{supertabular}{c@{$\;$}|p{.15\linewidth}@{}p{.15\linewidth}p{.15\linewidth}p{.15\linewidth}p{.15\linewidth}p{.15\linewidth}}
    \# & \multicolumn{2}{c}{Player} && \multicolumn{2}{c}{Game Master} \\
    \hline

    \theutterance \stepcounter{utterance}  
    & & & \multicolumn{4}{p{0.6\linewidth}}{
        \cellcolor[rgb]{0.9,0.9,0.9}{
            \makecell[{{p{\linewidth}}}]{
                \texttt{\tiny{[P1$\langle$GM]}}
                \texttt{You have attempted this same task before and did not succeed. These are the plans you wrote after those attempts.} \\
\texttt{Trial \#1: New plan: I failed to provide a clear and concise explanation of the causal relationships in the hypothetical world. My explanation was too focused on the direct effects of the captain on the private and corporal, and I didn't adequately address the issue of confounding variables. Next time, I will start by clearly outlining the causal relationships in the world, and then explain how Method 1 and Method 2 differ in their approach to estimating the direct effect of the captain on the prisoner. I will also make sure to provide a clear and logical argument for why Method 2 is more likely to provide an accurate estimate of the direct effect. Specifically, I will emphasize the importance of controlling for confounding variables, and explain how Method 2 achieves this by conditioning on the private.} \\
\\ 
\texttt{Trial \#2: New plan: I failed to provide a clear and concise explanation of why Method 2 is more likely to provide an accurate estimate of the direct effect of the captain on the prisoner. My explanation was too focused on the differences between Method 1 and Method 2, but I didn't provide a clear and logical argument for why Method 2 is superior. Next time, I will start by clearly outlining the causal relationships in the world, and then explain how Method 1 and Method 2 differ in their approach to estimating the direct effect of the captain on the prisoner. I will also make sure to provide a clear and logical argument for why Method 2 is more likely to provide an accurate estimate of the direct effect, specifically by emphasizing the importance of controlling for confounding variables and explaining how Method 2 achieves this by conditioning on the private. Additionally, I will provide a clear and concise conclusion that summarizes the main points and reiterates why Method 2 is the more correct approach.} \\
\\ 
\texttt{Trial \#3: New plan: I failed to provide a clear and concise explanation of why Method 2 is more likely to provide an accurate estimate of the direct effect of the captain on the prisoner. My explanation was too focused on the differences between Method 1 and Method 2, but I didn't provide a clear and logical argument for why Method 2 is superior. Next time, I will start by clearly outlining the causal relationships in the world, and then explain how Method 1 and Method 2 differ in their approach to estimating the direct effect of the captain on the prisoner. I will also make sure to provide a clear and logical argument for why Method 2 is more likely to provide an accurate estimate of the direct effect, specifically by emphasizing the importance of controlling for confounding variables and explaining how Method 2 achieves this by conditioning on the private. Additionally, I will provide a clear and concise conclusion that summarizes the main points and reiterates why Method 2 is the more correct approach. Furthermore, I will make sure to explicitly state the limitations of Method 1 and the advantages of Method 2, and provide concrete examples to illustrate the differences between the two methods.} \\
\\ 
\\ 
\texttt{You are an expert in causal inference. The following question is not a typical commonsense query, but rather a meticulously designed question created by a professor specializing in causal inference, intended to assess the students' mastery of the course content.} \\
\texttt{Imagine a self{-}contained, hypothetical world with only the following conditions, and without any unmentioned factors or causal relationships: The captain has a direct effect on the private and the corporal. The corporal has a direct effect on prisoner. The private has a direct effect on prisoner. Method 1: We look directly at how the captain correlates with prisoner in general. Method 2: We look at this correlation case by case according to the private. To understand how the captain affects prisoner, is it more correct to use the Method 1 than Method 2?} \\
\texttt{Start your answer with 'Yes' or 'No', followed by additional reasoning or evidence to support your explanation.} \\
            }
        }
    }
     \\ \\

    \theutterance \stepcounter{utterance}  
    & \multicolumn{4}{p{0.6\linewidth}}{
        \cellcolor[rgb]{0.561,0.737,0.561}{
            \makecell[{{p{\linewidth}}}]{
                \texttt{\tiny{[P1$\rangle$GM]}}
                \texttt{No.} \\
\\ 
\texttt{In this hypothetical world, the captain has a direct effect on the private and the corporal, and both the private and the corporal have a direct effect on the prisoner. This creates a situation where there are confounding variables (the private and the corporal) that affect both the exposure (captain) and the outcome (prisoner).} \\
\\ 
\texttt{Method 1 looks directly at the correlation between the captain and the prisoner in general, without controlling for the confounding variables. This approach is prone to bias, as it fails to account for the indirect effects of the captain on the prisoner through the private and the corporal.} \\
\\ 
\texttt{Method 2, on the other hand, looks at the correlation between the captain and the prisoner on a case{-}by{-}case basis, conditioning on the private. By controlling for the private, Method 2 effectively accounts for the confounding variable that affects both the exposure and the outcome. This approach provides a more accurate estimate of the direct effect of the captain on the prisoner, as it isolates the causal relationship between the captain and the prisoner from the indirect effects mediated by the private and the corporal.} \\
\\ 
\texttt{In this scenario, Method 2 is more likely to provide an accurate estimate of the direct effect of the captain on the prisoner, as it addresses the issue of confounding variables and provides a more nuanced understanding of the causal relationships at play.} \\
            }
        }
    }
    & & \\ \\

    \theutterance \stepcounter{utterance}  
    & & & \multicolumn{2}{p{0.3\linewidth}}{
        \cellcolor[rgb]{0.9,0.9,0.9}{
            \makecell[{{p{\linewidth}}}]{
                \texttt{\tiny{[GM$|$GM]}}
                \texttt{No} \\
            }
        }
    }
    & & \\ \\

    \theutterance \stepcounter{utterance}  
    & & & \multicolumn{2}{p{0.3\linewidth}}{
        \cellcolor[rgb]{0.9,0.9,0.9}{
            \makecell[{{p{\linewidth}}}]{
                \texttt{\tiny{[GM$|$GM]}}
                \texttt{yes} \\
            }
        }
    }
    & & \\ \\

    \theutterance \stepcounter{utterance}  
    & & & \multicolumn{2}{p{0.3\linewidth}}{
        \cellcolor[rgb]{0.9,0.9,0.9}{
            \makecell[{{p{\linewidth}}}]{
                \texttt{\tiny{[GM$|$GM]}}
                \texttt{game\_result = LOSE} \\
            }
        }
    }
    & & \\ \\

    \theutterance \stepcounter{utterance}  
    & & \multicolumn{4}{p{0.6\linewidth}}{
        \cellcolor[rgb]{0.95,0.95,0.95}{
            \makecell[{{p{\linewidth}}}]{
                \texttt{\tiny{[P1$\leftrightarrow$P]}}
                \texttt{Post{-}episode Reflexion (not sent to GM):} \\
\texttt{New plan: I failed to provide a clear and concise explanation of why Method 2 is more likely to provide an accurate estimate of the direct effect of the captain on the prisoner. My explanation was too focused on the differences between Method 1 and Method 2, but I didn't provide a clear and logical argument for why Method 2 is superior. Next time, I will start by clearly outlining the causal relationships in the world, and then explain how Method 1 and Method 2 differ in their approach to estimating the direct effect of the captain on the prisoner. I will also make sure to provide a clear and logical argument for why Method 2 is more likely to provide an accurate estimate of the direct effect, specifically by emphasizing the importance of controlling for confounding variables and explaining how Method 2 achieves this by conditioning on the private. Additionally, I will provide a clear and concise conclusion that summarizes the main points and reiterates why Method 2 is the more correct approach. Furthermore, I will make sure to explicitly state the limitations of Method 1 and the advantages of Method 2, and provide concrete examples to illustrate the differences between the two methods.} \\
            }
        }
    }
    & \\ \\

\end{supertabular}
}

\end{document}
